\chapter{Introducción}\label{cap:introduccion}

\section{Contexto}


\section{Motivación}


\section{Estado del Arte}

En el presente apartado se ofrece una visión general del dominio del problema, abarcando tanto los fundamentos eléctricos como las técnicas de aprendizaje automático relevantes para la comprensión del objetivo del presente proyecto.

\subsection{Redes de Distribución}

\subsection{Técnicas de Aprendizaje Automático}

El Aprendizaje Automático (Machine Learning) es una rama de la Inteligencia Artifical cuyo objetivo es el estudio y desarrollo de sistemas capaces de adaptar su comportamiento a partir de la experiencia. Para construir dicha experiencia, el sistema aplica un algoritmo de aprendizaje sobre un conjunto de ejemplos de entrenamiento, dotando al sistema de la capacidad de aprender y predecir resultados sin necesidad de una programación explícita. 

Existen diversos tipos de algoritmos de Aprendizaje Automático, que pueden clasificarse en general en cuatro categorías principales:
\begin{itemize}
    \item \textbf{Aprendizaje supervisado:} El sistema se entrena con un conjunto de datos etiquetados. Cada ejemplo se describe mediante un conjunto de atributos de entrada o características y una salida objetivo conocida, que define la salida pretendida por el sistema. El objetivo del modelo es aprender la relación entre las entradas y las salidas para predecir correctamente la salida de nuevos ejemplos.
    \item \textbf{Aprendizaje no supervisado:} El sistema se entrena con un conjunto de datos sin etiquetar. De cada ejemplo se desconoce la salida esperada, por lo tanto su objetivo es identificar patrones o estructuras en los datos que ayuden a organizarlos.
    \item \textbf{Aprendizaje semisupervisado:} El sistema se entrena con un conjunto de datos en el que solo una parte está etiquetada. De esta forma se aprovecha tanto la información supervisada como la no supervisada para mejorar el rendimiento del modelo.
    \item \textbf{Aprendizaje por refuerzo:} El sistema aprende mediante la interración con el entorno, tomando decisiones en una secuencia de estados y recibiendo recompensas o castigos según el resultado de sus acciones. El objetivo es elegir la mejor secuencia de acciones (política óptima) que maximiza la recompensa acumulada.
\end{itemize}

En este proyecto nos centraremos en algoritmos de aprendizaje supervisado, que es el enfoque más adecuado para el objetivo de este, ya que se dispone de un dataset con el estado de la red en cada momento y la salida objetivo correspondiente. Estos algoritmos se pueden dividir en clasificación, cuando la variable objetivo referencia a tipo discreto y en regresión cuando la variable objetivo referencia a un valor continuo.

\subsubsection{Random Forest}

\subsubsection{XGBoost}

\subsubsection{Redes Neuronales}

\subsubsection{Otras Opciones}

\section{Objetivos}

El objetivo del presente Trabajo Fin de Grado consiste en desarrollar y evaluar diferentes modelos de Inteligenia Artificial mediante aprendizaje supervisado, para predecir valores de referencia (set points) que optimicen el funcionamiento de redes eléctricas. El propósito de estos modelos será predecir de manera efectiva y precisa estos valores de referencia, reduciendo la necesidad de disponer de información detallada sobre la topología y parámetros eléctricos de la red, requeridos por métodos tradicionales de optimización como el algoritmo Optimal Power Flow. Con este fin el entramiento de estos modelos se realizarán a partir de grandes volúmenes de datos medidos por sensores instalados en la red.

Para alcanzar este objetivo general, se plantean los siguientes objetivos específicos:
\begin{enumerate}
    \item \textbf{TODO: Lista de sub-objetivos a realizar} %TODO: LISTADO DE OBJETIVOS
\end{enumerate}

%TODO: INCLUIR O NO? \section{Estructura}